\bookchapter{第十五章：电池科学与续航优化}

\chapterintro{电池就是你手机的“能量心脏”。这章带你搞懂电池的小脾气，学会科学的充电和续航技巧，彻底告别电量焦虑。}

\booksection{一、电池是怎么回事？循环与健康度}

手机电池（锂离子或锂聚合物）像个能量海绵：充电就是吸水（存电），放电就是挤水（用电）。有几个特点值得记住：

\begin{itemize}
  \item \textbf{随用随充}：和旧电池不同，它没有记忆效应，不用非得用完再充。
  \item \textbf{怕热}：高温是电池头号杀手。
  \item \textbf{会老化}：用久了，最大容量就会下降，这是正常现象。
\end{itemize}

\booksubsection{核心指标：循环次数与健康度}

\begin{itemize}
  \item 一次循环不是充一次电。比如今天用60\%充满，明天用40\%充满，总消耗100\%电量才算一次循环。
  \item 电池健康度就是当前最大容量占全新时的百分比。低于80\%时，你会明显感觉续航变短。
  \item \textbf{怎么看？}iPhone：设置 > 电池 > 电池健康与充电；安卓：一般去设置里搜电池健康。
\end{itemize}

\booksection{二、充电玄学？快充、无线充与黄金法则}

\begin{itemize}
  \item \textbf{快充伤电池？}用原装或认证的快充配件，伤害很小。手机和充电头有温控，过热会自动降功率。
  \item \textbf{无线充电}：方便但效率较低，发热更明显。适合整夜充电。
  \item \textbf{黄金养生法则}：尽量让电量保持在20\%-80\%之间，别老让它饿着（0\%）或撑着（100\%）。
\end{itemize}

\begin{tcbinfo}{iPhone"80\%充电限制"功能——把黄金法则变成系统设置}
iPhone 15/16系列（iOS 17以上）已经把“20\%-80\%法则”做成了系统功能。开启后，手机充到80\%会自动停止充电，从而大幅延缓电池老化。设置路径：

\texttt{设置 > 电池 > 电池健康与充电 > 充电优化 > 80\%限制}

三个选项的含义：(1)\textbf{不限}——充到100\%；(2)\textbf{优化电池充电}——学习你的作息，只在凌晨最后一段延迟到100\%；(3)\textbf{80\%限制}——直接停在80\%，最护电池。长期在家/办公室有充电条件的人，强烈建议选80\%限制；经常出差、需要满电出门的人，可以选优化电池充电。

\end{tcbinfo}

\begin{tcbinfo}{GaN充电器——出差党的好搭档}
你可能注意到现在的快充头越做越小，背后功臣是\textbf{GaN（氮化镓）}这种新型半导体材料。相比传统硅基充电器，GaN充电器有几个明显优势：

\begin{itemize}
  \item \textbf{体积小}：同等功率下，GaN充电器可以做到传统硅基的1/2甚至1/3大小，65W功率的GaN头比苹果老的5W头还小。
  \item \textbf{效率高、发热低}：氮化镓的开关损耗远低于硅，能量转换效率可达95\%以上，充电时不容易烫手。
  \item \textbf{适合出差}：一个65W或100W的GaN多口充电器，能同时给手机、平板、笔记本充电，比带三块"砖头"轻便得多。
  \item \textbf{选购建议}：认准支持PD协议的GaN充电器（65W起步可充笔电），知名品牌如倍思、安克、努比亚、小米等均可。注意核对数据线是否匹配（高功率需用E-Marker芯片的Type-C线）。
\end{itemize}

\end{tcbinfo}

\booksection{三、快充协议全解析：为什么你的充电器充不快？}

你可能遇到过这种情况：明明手机支持快充，换了充电器却变得很慢。这是因为快充不是简单的“功率大就行”，而是需要充电头、数据线、手机三方“说同一种语言”——这就是快充协议。

\booksubsection{快充的基本原理}

充电功率 = 电压 × 电流（P = U × I）。要提升功率，要么提高电压，要么提高电流，或者两者都提高。不同协议选择了不同的技术路线。

快充的工作流程：

\begin{enumerate}
  \item \textbf{握手阶段}：手机和充电头通过数据线通信，交换各自支持的协议和功率信息。
  \item \textbf{协商阶段}：双方确定一个共同支持的最高功率档位。
  \item \textbf{充电阶段}：充电头按协商好的电压电流输出，手机端进行电源管理。
  \item \textbf{调整阶段}：随着电量增加，手机会逐步降低充电功率，保护电池。
\end{enumerate}

\booksubsection{主流快充协议详解}

\textbf{USB PD（Power Delivery）}——通用协议之王

\begin{itemize}
  \item \textbf{特点}：USB-IF组织制定的官方标准，覆盖范围最广。PD 3.0支持最高100W（20V/5A），PD 3.1扩展到240W（48V/5A）。
  \item \textbf{电压档位}：5V、9V、15V、20V，PD 3.1还支持28V、36V、48V。
  \item \textbf{兼容性}：iPhone全系、谷歌Pixel、部分安卓旗舰、笔记本电脑、Switch游戏机都支持PD协议。
  \item \textbf{优势}：通用性强，一个充电头可以给多种设备充电；支持PPS（可编程电源），可以精细调节电压电流。
  \item \textbf{劣势}：需要Type-C接口和专用芯片，成本较高。
\end{itemize}

\textbf{高通QC（Quick Charge）}——安卓阵营的元老

\begin{itemize}
  \item \textbf{QC 3.0}：支持3.6V-20V以0.2V为步进的可调电压，最高18W。
  \item \textbf{QC 4.0/4+}：兼容USB PD协议，支持最高27W，加入双路充电技术。
  \item \textbf{QC 5.0}：支持最高100W+，0-50\%仅需5分钟，大幅改善发热。
  \item \textbf{兼容性}：使用高通骁龙处理器的手机大多支持，但各品牌通常会优先使用自己的私有协议。
\end{itemize}

\textbf{OPPO VOOC/一加WARP/真我Dart}——低压大电流的代表

\begin{itemize}
  \item \textbf{技术路线}：低压大电流（5V/4A起步），发热更小，效率更高。
  \item \textbf{VOOC 1.0}：20W（5V/4A），2014年发布，开创快充时代。
  \item \textbf{SuperVOOC}：50W-65W，采用双电芯串联设计。
  \item \textbf{SuperVOOC 2.0}：65W，30分钟充满4000mAh电池。
  \item \textbf{150W/240W超级闪充}：目前最快的量产快充技术。
  \item \textbf{特点}：需要专用充电头和数据线（有额外触点），不兼容其他协议。
\end{itemize}

\textbf{华为SuperCharge}——高压快充的演进

\begin{itemize}
  \item \textbf{FCP（Fast Charge Protocol）}：9V/2A，18W，兼容QC 2.0。
  \item \textbf{SCP（Super Charge Protocol）}：支持更高功率，从22.5W到66W再到100W。
  \item \textbf{技术特点}：采用电荷泵技术，将高压转换为低压大电流输入电池，效率高达98\%。
  \item \textbf{兼容性}：华为/荣耀手机专用，部分充电器支持PD协议。
\end{itemize}

\textbf{小米澎湃快充}——自研协议的崛起

\begin{itemize}
  \item \textbf{澎湃P1芯片}：小米自研的快充管理芯片，支持120W有线快充。
  \item \textbf{澎湃G1芯片}：电池管理芯片，与P1配合实现更高效的快充。
  \item \textbf{技术特点}：采用双电芯串联+电荷泵方案，支持多种协议兼容。
\end{itemize}

\textbf{苹果快充}——保守但稳定

\begin{itemize}
  \item \textbf{协议}：仅支持USB PD协议。
  \item \textbf{功率}：iPhone 8及以后支持快充，最高约27W（iPhone 15 Pro Max可达27W+）。
  \item \textbf{特点}：功率相对保守，但充电曲线优化好，发热控制出色。
  \item \textbf{配件要求}：需要USB-C to Lightning线（或USB-C to USB-C线）和PD充电头。
\end{itemize}

\booksubsection{快充协议对比表}

\begin{center}
\begin{tabularx}{\linewidth}{|X|X|X|X|}
\hline
\textbf{协议} & \textbf{最高功率} & \textbf{技术路线} & \textbf{兼容性} \tabularnewline
\hline
USB PD 3.1 & 240W & 高压/多档位 & 最广泛 \tabularnewline
QC 5.0 & 100W+ & 高压可调 & 高通机型 \tabularnewline
VOOC/Dart & 240W & 低压大电流 & OPPO系专用 \tabularnewline
SuperCharge & 100W & 电荷泵 & 华为/荣耀 \tabularnewline
苹果快充 & \textasciitilde{}27W & PD协议 & iPhone/iPad \tabularnewline
\hline
\end{tabularx}
\end{center}

\booksubsection{无线充电协议}

无线充电也有自己的协议标准：

\begin{itemize}
  \item \textbf{Qi标准}：WPC（无线充电联盟）制定，最通用的标准。基础Qi支持5W-15W。
  \item \textbf{Qi2标准}：2023年发布，引入磁吸定位（类似苹果MagSafe），支持15W+。
  \item \textbf{私有协议}：各品牌有自己的高速无线充电，如小米50W无线充、华为50W无线充等，需要专用充电器。
\end{itemize}

\booksubsection{快充选购指南}

\begin{itemize}
  \item \textbf{看协议匹配}：购买充电器时，确认它支持你手机的快充协议。第三方充电器通常会标注支持的协议。
  \item \textbf{看数据线}：大功率快充需要专用数据线。普通数据线可能只支持60W以下。
  \item \textbf{看认证}：苹果用户建议选择MFi认证的数据线；安卓用户选择品牌原装或知名第三方。
  \item \textbf{看发热}：快充时手机发热是正常的，但如果烫手就不正常了，建议暂停充电或更换充电器。
\end{itemize}

\begin{tcbwarn}{快充安全提示}
\begin{itemize}
  \item 避免边充边玩大型游戏，会加剧发热。
  \item 不要使用劣质充电器和数据线，可能导致安全隐患。
  \item 充电时取掉厚重的手机壳，帮助散热。
  \item 如果充电器或手机异常发热、有异味，立即停止使用。
\end{itemize}

\end{tcbwarn}

\booksection{四、电池新技术：硅碳负极电池}

如果你留意2024-2026年的旗舰手机发布会，“硅碳负极”这个词几乎成了高端电池的代名词。这是锂电池负极材料的一次重要升级，原理其实不复杂：

\begin{itemize}
  \item \textbf{传统石墨负极}：手机电池一直用的是石墨做负极，技术成熟、稳定，但能量密度已接近理论上限（约372mAh/g）。
  \item \textbf{掺入硅材料}：硅的理论比容量是石墨的10倍左右（约4200mAh/g）。把少量硅掺入石墨里做成"硅碳复合负极"，就能显著提升整块电池的能量密度。
  \item \textbf{能量密度提升}：同等体积下，硅碳负极电池的容量可以提升约\textbf{15\%-20\%}。这正是2024-2026年旗舰机能在7-9mm机身里塞进5500-6000mAh大电池的关键原因。
\end{itemize}

\begin{tcbinfo}{哪些机型用上了？}
\begin{itemize}
  \item \textbf{荣耀}：Magic系列率先商用硅碳负极，Magic6 Pro、Magic7系列、Magic V3折叠屏均采用。
  \item \textbf{小米}：小米14 Ultra、小米15系列搭载新一代硅碳负极电池。
  \item \textbf{vivo}：X100 Ultra、X200 Pro系列使用高密度硅碳负极方案。
  \item \textbf{华为/OPPO等}：2025-2026年旗舰普遍跟进，硅碳负极正在从"黑科技"变成高端机的标配。
\end{itemize}

\end{tcbinfo}

\textbf{对用户的实际影响}：硅碳负极带来的最大好处是“同体积电池容量更大”——手机可以做得更薄，同时续航不缩水；或者机身不变厚，但电池从4500mAh涨到5500mAh。换句话说，你不用为续航付“厚重”的代价了。需要注意的是，硅碳负极电池在低温和满充管理上仍需配合快充芯片做优化，这也是厂商建议长期保持80\%充电的原因之一。

\booksection{五、这才是真正的省电技巧}

\begin{itemize}
  \item \textbf{屏幕（耗电王）}：调低亮度、缩短自动锁屏时间。OLED屏手机开深色模式能省不少电。
  \item \textbf{信号差}：在电梯、地下室，手机会拼命搜信号，耗电极快。
  \item \textbf{5G}：如果信号不好，手动切回4G，续航会好很多。
  \item \textbf{定位服务}：非导航类App，定位权限设为"仅使用时允许"。
  \item \textbf{限制后台活动}：iPhone在设置 > 通用 > 后台App刷新；安卓在应用电池设置里。
\end{itemize}

\begin{tcbpractice}{来检查一下：给你的电池做个体检}
\begin{enumerate}
  \item 查看健康度：找到并记下你的电池最大容量。
  \item 分析耗电详情：看电池设置里的图表和App列表，找出耗电元凶。
  \item 制定对策：限制后台耗电异常的App。
\end{enumerate}

\end{tcbpractice}
